\documentclass[10pt,a4paper]{amsart}
\usepackage[margin=19mm]{geometry}
\usepackage{amsmath,amssymb,mathtools}
\usepackage[scr=rsfs,cal=euler]{mathalfa}
\usepackage{xfrac}
\usepackage[unicode,hidelinks,bookmarks=false]{hyperref}
\newcommand{\ie}{i.e.\ }
\newcommand{\cf}{cf.\ }
\newcommand{\resp}{respectively\ }
\newcommand{\Subsection}{\subsection}
\providecommand{\nobreakdash}{\nobreak\hbox{-}}
\setlength{\emergencystretch}{2em}
\multlinegap0pt
\usepackage{amssymb}
\usepackage{mathtools}
\usepackage{xcolor}
\usepackage[shortlabels]{enumitem}
\newtheorem{lemma}{Lemma}
\newtheorem{theorem}{Theorem}
\newcommand{\Connection}{A}
\newcommand{\Op}{\mathcal{O}}              % Matérn operator
\newcommand{\Z}{\mathbb{Z}}
\newcommand{\inj}{\operatorname{inj}}      % injectivity radius
\newcommand{\sph}[1]{S^{#1}}
\title{Yang--Mills--Higgs Classification: \\ A Geometric Theory of Binary Labels on Non-Contractible Spaces}
\author{Catalin Vasii}
\date{Source: arXiv:2607.00999v2 (CC BY 4.0)}
\begin{document}
\maketitle

\noindent\emph{Source and licence.} Yang--Mills--Higgs Classification: \\ A Geometric Theory of Binary Labels on Non-Contractible Spaces. Version arXiv:2607.00999v2 (CC BY 4.0), distributed by arXiv under the Creative Commons Attribution 4.0 International licence, http://creativecommons.org/licenses/by/4.0/, as recorded in the arXiv licence metadata for this exact version and shown on the abstract page https://arxiv.org/abs/2607.00999v2. Full licence notice: \texttt{licenses/arxiv-2607.00999-CC-BY-4.0.txt}.\par\smallskip

\noindent\textit{Editorial scope.} Counted unit: the article's theorem thm:proximity-scaling-2pt (source0.tex lines 5348--5378). The article's complete original proof of it is delivered with it (source0.tex lines 5380--5415, one \texttt{proof} environment of 36 lines). No mathematics is written, repaired or re-derived: the statement and the proof are the article's own lines, in the article's own notation, and no retained line is substituted or re-flowed.  Retained uncounted as the source-local context the proof needs: lines 5189--5192; lines 5241--5284.  Nothing was omitted from any retained span, so the package carries no omission record.  No retained line contains a citation, so the wrapper carries no bibliography and no bibliography entry.  No retained line contains a figure environment, an image-inclusion command or drawing code, so no image is delivered and the package declares no asset dependency.  Editorial formatting only, in addition: an AMS \texttt{amsart} preamble that restates the article's own declarations and macros used by the retained lines, namely the environments \texttt{lemma}, \texttt{theorem} on separate counters, together with the standard packages those lines need.  The retained environments print in this document as Lemma 1, Theorem 1 in order of appearance, whereas the article prints the same items with the numbering of its own full document.  The complete native source of the article as distributed in the arXiv e-print for this exact version is bundled unchanged as \texttt{sources/204202/source0.tex}.
\begin{equation}\label{eq:cov-capacitance}
   K\,\alpha \;=\; y,
   \qquad K_{ij} := G_A(x_i,x_j),\quad y=(y_1,\dots,y_N).
\end{equation}

\begin{lemma}[Local asymptotic of $G_A$]
\label{lem:kernel-asymptotic}
Let $(M,g)$ be closed and Riemannian, $L_\xi$ a flat line bundle,
$\Op_A=(D_A^*D_A+\kappa^2)^\nu$ with $\nu>n/2,\ \kappa>0$
(the hypothesis $\nu>n/2$ ensures that $G_A(x,x)$ is finite on the
diagonal --- a Sobolev/parametrix bound that the lemma uses). Then
as $d\to 0^+$,
\[
   G_A(0)-G_A(d) \;=\;
   \begin{cases}
      c_{n,\nu}\,d^{2\nu-n} + o(d^{2\nu-n})
      & \text{if } 2\nu-n\in(0,2)\text{ and }2\nu-n\notin 2\Z, \\[3pt]
      c_{n,\nu}'\,d^{2\nu-n}\log(1/d) + O(d^{2\nu-n})
      & \text{if } 2\nu-n\in 2\Z_{>0}\cap(0,2], \\[3pt]
      C\,d^2 + o(d^2) & \text{if } 2\nu-n\geq 2 \text{ generic,}
   \end{cases}
\]
where $c_{n,\nu}, c_{n,\nu}'>0$ are dimensional constants depending
only on $n$ and $\nu$ (with
$c_{n,\nu}=-\,\Gamma(n/2-\nu)/(2^{2\nu}\pi^{n/2}\Gamma(\nu))$
for $2\nu-n\notin\Z$, and analogous expressions for the log branch;
the explicit minus sign is required because $\Gamma$ is evaluated at
the negative argument $n/2-\nu<0$ when $2\nu-n>0$, so that
$c_{n,\nu}>0$ as the statement requires --- for $n=1,\nu=1$ it gives
$c_{1,1}=1/2$, matching the exact $\sph{1}$ value
$G(0)-G(d)=d/2$).
Logarithmic factors arise precisely at \emph{even} positive integer
values of $2\nu-n$, because at those values the Bessel index
$\nu-n/2$ is a non-negative integer and the small-argument expansion
of $K_{\nu-n/2}$ acquires a logarithmic term; at odd or non-integer
values of $2\nu-n$, the Bessel index is a half-integer or
non-integer and $K_{\nu-n/2}$ has a pure-power expansion (in the
half-integer case it is even elementary, with no log).

The bundle structure of $L_\xi$ contributes only at subleading
order: for $d$ smaller than the injectivity radius and any chart
trivialising $L_\xi$, $G_A(x_+,x_-)$ equals the trivial-bundle
scalar kernel up to corrections one order down from the leading
$d^{2\nu-n}$ term, i.e.\ at order $d^{2\nu-n+1}$ or $d^{2\nu-n+2}$
(depending on the parity of $2\nu-n$ and the geometry of $(M,g)$).
This subleading bound uses crucially the flatness of $\Connection$:
the argument fails for curved connections, where the symbol-level
gauge-conjugation of Step~1 below breaks down.
\end{lemma}

\begin{theorem}[Matter-sector proximity scaling, two-point case]
\label{thm:proximity-scaling-2pt}
Let $(M,g)$ be a closed Riemannian $n$-manifold, $L_\xi$ a flat
$\Z_2$-bundle, $\Op_A=(D_A^*D_A+\kappa^2)^\nu$ with $\nu>n/2,\
\kappa>0$ (the $\nu>n/2$ hypothesis ensures $G_A(0)$ is finite, as
required for the energy formula below). Take $N=2$ labelled points
$x_+,x_-$ with opposite labels $y_\pm=\pm 1$, separated by geodesic
distance $d<\inj(M,g)$, and assume the variational selector picks
the bundle $L_\xi$ such that the parallel section reverses sign
between $x_+$ and $x_-$ along the geodesic. Then
\[
   E_{\min} \;=\; \frac{2}{G_A(0)-G_A(d)},
\]
and the small-$d$ asymptotic of $E_{\min}$ is determined by
Lemma~\ref{lem:kernel-asymptotic}:
\[
   E_{\min} \;\sim\;
   \begin{cases}
      \dfrac{2}{c_{n,\nu}}\,d^{-(2\nu-n)}
      & \text{if } 2\nu-n\in(0,2)\setminus 2\Z, \\[6pt]
      \dfrac{2}{c_{n,\nu}'}\,\dfrac{d^{-(2\nu-n)}}{\log(1/d)}
      & \text{if } 2\nu-n\in 2\Z_{>0}\cap(0,2], \\[8pt]
      \dfrac{2}{C}\,d^{-2}
      & \text{if } 2\nu-n>2,
   \end{cases}
\]
as $d\to 0^+$, with constants from
Lemma~\ref{lem:kernel-asymptotic}. In particular, the scaling
exponent saturates at $\min(2\nu-n,\,2)$ once $\nu$ is large enough,
with logarithmic factors at the boundary $2\nu-n\in 2\Z_{>0}$.
\end{theorem}

\begin{proof}
Work in the local trivialisation of $L_\xi$ on a neighbourhood
covering the geodesic from $x_+$ to $x_-$ (possible since
$d<\inj(M,g)$); the flat connection $\Connection$ becomes a locally
exact 1-form $df$ on this chart, and the covariant kernel
$G_A(x_i,x_j)$ for $x_i,x_j$ in the chart equals the scalar Mat\'ern
kernel $G(d(x_i,x_j))$ on $(M,g)$ up to gauge-conjugation by $e^f$
(which preserves the kernel diagonal $G_A(0)=G(0)$ and the off-diagonal
$G_A(x_+,x_-)=G(d)$ for points in the same chart). The data
conditions in this trivialisation are $\phi^*(x_+)=+1$ and
$\phi^*(x_-)=-1$: the sign-reversing parallel transport that defines
$L_\xi$ globally is encoded in the bundle's transition function on
the overlap with a second chart, not in the values on a single chart.
The capacitance system \eqref{eq:cov-capacitance} for $N=2$ in this
trivialisation is therefore
\[
   \begin{pmatrix} G(0) & G(d) \\ G(d) & G(0) \end{pmatrix}
   \begin{pmatrix}\alpha_+\\\alpha_-\end{pmatrix}
   \;=\; \begin{pmatrix}+1\\-1\end{pmatrix},
\]
giving $\alpha_\pm=\pm 1/(G(0)-G(d))$. The minimum energy is
\[
   E_{\min} \;=\; \alpha^\top K\alpha
   \;=\; y^\top\alpha
   \;=\; \frac{2}{G(0)-G(d)}.
\]
The bundle's M\"obius character (selected by the variational
selector as carrying lower energy than the trivial bundle in this
configuration) does not affect the capacitance solve in this single
trivialising chart; its effect is to select \emph{which} bundle's
Green's kernel is the correct one, with the M\"obius bundle's
chart-local kernel agreeing with the scalar Mat\'ern kernel of
$(M,g)$ to leading order by Lemma~\ref{lem:kernel-asymptotic}.
Applying that lemma to the denominator yields the stated $d\to 0^+$
asymptotic.
\end{proof}

\end{document}
